% INTRODUÇÃO: substitua as orientações pelo texto da sua pesquisa.
% Preserve o identificador do \label se ele for citado em outras seções.
\section{Introdução}
\label{sec:introducao}

Apresente o tema e delimite o problema que será investigado. Explique por que
o estudo é relevante, formule a pergunta de pesquisa e indique o objetivo
geral e os objetivos específicos. Situe o leitor no contexto do trabalho.

Este modelo contém textos, imagens e referências de demonstração. Substitua-os
pelos materiais da sua pesquisa antes da entrega. A aparência do documento é
controlada pelos arquivos de configuração; não ajuste cada parágrafo manualmente.

% Exemplo de citação: a chave identifica uma entrada em bibliografia.bib.
Uma citação entre parênteses é produzida assim: \cite{SilvaNeto2019Credibility}.
Este exemplo demonstra o comando, não estabelece uma relação entre a obra
citada e o tema que você escolher.

% Tabela com largura ajustável: X ocupa o espaço que sobra na linha.
A \autoref{tab:identificador} demonstra uma tabela com quatro colunas.

\begin{table}[htb]
  \centering
  \ABNTEXfontereduzida
  \caption{Exemplo de organização de informações}
  \label{tab:identificador}
  \begin{tabularx}{\linewidth}{lllX}
    \toprule
    \textbf{Item} & \textbf{Categoria} & \textbf{Etapa} & \textbf{Descrição} \\
    \midrule
    1 & Exemplo & Inicial & Substitua este conteúdo pelos dados do trabalho.
    A última coluna ajusta sua largura automaticamente. \\
    \bottomrule
  \end{tabularx}
  \fonte{Elaboração própria; dados fictícios para demonstração.}
\end{table}

% Atualize este parágrafo caso acrescente ou remova seções.
O trabalho está organizado em introdução (Seção~\ref{sec:introducao}),
fundamentação teórica (Seção~\ref{sec:fundamentacao}), metodologia
(Seção~\ref{sec:metodologia}), resultados e discussão
(Seção~\ref{sec:resultados-discussoes}) e conclusão
(Seção~\ref{sec:conclusao}).
